In the following, we want to use the Banach contraction principle to prove an implicit function theorem for functions between Banach spaces. In order to do so, we first need to define the concept of a derivative.

\begin{Definition}
\label{def:frechet-derivative}
Let $X$ and $Y$ be Banach spaces and let $x_0$ be a point in $X$. Let $F$ be a mapping from a neighborhood of $x_0$ into $Y$. Then $F$ is called \textbf{differentiable} at $x_0$ if there exists a linear operator $A\in\mathcal{L}(X,Y)$ with the property that
\begin{equation}\label{eq:frechet-differentiability-linear-remainder-expansion}
F(x)=F(x_0)+A(x-x_0)+G(x),
\end{equation}
where
\begin{equation}\label{eq:frechet-differentiability-vanishing-relative-remainder}
\lim_{x\to x_0}\|G(x)\|_Y/\|x-x_0\|_X=0.
\end{equation}
If such an $A$ exists, we call it the (Fréchet) derivative of $F$ at $x_0$.
\end{Definition}

Naturally, we shall call a function differentiable if it is differentiable at all points of its domain. We shall use familiar notations from calculus, such as $F^{\prime}(x_0)$ or $DF(x_0)$ to denote Fréchet derivatives. We call a function continuously differentiable, if $F^{\prime}(x)$ (as an element of $\mathcal{L}(X,Y)$) depends continuously on $x$. Clearly, we can use the definition recursively to define higher-order derivatives. For example, $F^{\prime\prime}(x)$, if it exists, is an element of $\mathcal{L}(X,\mathcal{L}(X,Y))$. Alternatively, we can regard $F^{\prime\prime}(x)$ as a bilinear mapping from $X$ to $Y$, see Problem~\ref{ex:continuous-bilinear-linear-operator-correspondence}. Many of the elementary properties of derivatives can be generalized to Fréchet derivatives in a straightforward manner; see, e.g., Problem~\ref{ex:frechet-derivative-chain-rule}.

\begin{Example}
\label{ex:pointwise-composition-frechet-derivative-example}
Let $f$ be a continuously differentiable function from $\mathbb{R}$ to $\mathbb{R}$ and let $X=Y=C([0,1])$. The mapping $f$ induces a mapping from $X$ to $Y$ by pointwise action: $\phi(t)\mapsto f(\phi(t))$. The mapping $f$ is differentiable as a mapping from $X$ to $Y$ and its derivative at $\phi$ is the linear mapping which takes the function $\psi(t)$ to $f^{\prime}(\phi(t))\psi(t)$. See Problem~\ref{ex:pointwise-composition-frechet-derivative}.
\end{Example}
